\section*{Useful links}

A small collection of websites that help when writing a report.
\begin{footnotesize}
    \begin{multicols}{2}
        \begin{itemize}
            \item \href{https://molview.org/}{MolView}: 2D and 3D images of chemical compounds;
            \item \href{https://py-chemist.com/mol_2_chemfig/home}{Mol2chemfig}: draws molecules and mechanisms and converts them into \LaTeX\ code for the \texttt{chemfig} package;
            \item \href{https://echa.europa.eu/regulations/clp/clp-pictograms}{ECHA -- CLP pictograms}: the European hazard pictograms for the labelling of chemicals;
            \item \href{https://echa.europa.eu/information-on-chemicals/cl-inventory-database}{ECHA -- C\&L Inventory}: classification and labelling of substances in the EU;
            \item \href{https://pubchem.ncbi.nlm.nih.gov/}{PubChem}: the world's largest collection of freely accessible chemical information: properties, safety and toxicity data, literature and patents;
            \item \href{https://www.chemspider.com/}{ChemSpider}: free chemical structure database with text and structure search over more than 100 million structures;
            \item \href{https://commonchemistry.cas.org/}{CAS Common Chemistry}: open resource with about \num{500000} substances from the CAS REGISTRY, including the common ones used in school and undergraduate chemistry;
            \item \href{https://www.ebi.ac.uk/chebi/}{ChEBI}: Chemical Entities of Biological Interest, a dictionary of small molecular entities;
            \item \href{https://www.ebi.ac.uk/chembl/}{ChEMBL}: curated database of bioactive molecules with drug-like properties;
            \item \href{https://go.drugbank.com/}{DrugBank}: drug interactions, pharmacology, chemical structures, targets and metabolism;
            \item \href{https://www.rcsb.org/}{RCSB PDB}: Protein Data Bank, the public archive of 3D structures of proteins and nucleic acids;
            \item \href{https://www.uniprot.org/}{UniProt}: protein sequences and functional information;
            \item \href{https://enviPath.org/}{enviPath}: biodegradation pathways of chemicals (successor of the EAWAG-BBD);
            \item \href{https://www.efsa.europa.eu/}{EFSA}: European Food Safety Authority, scientific risk assessments and data;
            \item \href{https://www.ars.usda.gov/}{ARS}: Agricultural Research Service of the US Department of Agriculture;
            \item \href{https://www.nist.gov/}{NIST}: National Institute of Standards and Technology; the \href{https://webbook.nist.gov/chemistry/}{NIST Chemistry WebBook} provides thermochemical and spectroscopic data;
            \item \href{https://www.rsc.org/merck-index}{The Merck Index Online}: authoritative reference on chemicals, drugs and biologicals;
            \item \href{https://www.webelements.com/}{WebElements}: the periodic table and the properties of the elements;
            \item \href{https://www.chm.bris.ac.uk/motm/motm.htm}{Molecule of the Month}: one of the longest-running chemistry websites, with a new molecule each month since 1996;
            \item \href{https://chemistry.stackexchange.com/}{Chemistry Stack Exchange}: questions and answers on chemistry;
            \item \href{https://tex.stackexchange.com/}{\TeX\ Stack Exchange}: questions and answers on \LaTeX;
            \item \href{https://ctan.org/pkg/mhchem}{mhchem}, \href{https://ctan.org/pkg/chemfig}{chemfig}, \href{https://ctan.org/pkg/siunitx}{siunitx}, \href{https://ctan.org/pkg/pgfplots}{pgfplots}: documentation of the main packages used by this template;
            \item \href{https://ctan.org/pkg/annotate-equations}{annotate-equations}: package to annotate the terms of an equation.
        \end{itemize}
    \end{multicols}
\end{footnotesize}
